% Recommended, but optional, packages for figures and better typesetting:
\usepackage{microtype}
\usepackage{graphicx}
\usepackage{booktabs} % for professional tables
% \usepackage{fontspec}
% to be able to draw some self-contained figs
\usepackage{tikz}
\usepackage{amsmath}
\usepackage{multirow}
% inlined bib file
% \usepackage{filecontents}
\usepackage{fontawesome}
\usepackage{subcaption}
\usepackage{xspace}
\usepackage{booktabs}
\usepackage{comment}
\usepackage{float}
\usepackage{caption}
\usepackage{lipsum}
\usepackage[]{hyperref}
\usepackage{comment}
\usepackage{xspace}
\usepackage{wasysym}
\usepackage{float}
\usepackage{stfloats}
\usepackage{placeins}
\usepackage{enumitem}
\usepackage{xcolor}
\usepackage{booktabs}
\PassOptionsToPackage{table}{xcolor}
\PassOptionsToPackage{x11names}{xcolor}
\usepackage{xcolor}
\usepackage{colortbl}
\usepackage{makecell}
\usepackage{nicematrix}
\usepackage[ruled,linesnumbered]{algorithm2e}
\usepackage[export]{adjustbox}
\usepackage{float}
\usepackage[noend]{algpseudocode}
\usepackage{tabularx} % In the preamble

% Markers
\let\Bbbk\relax
\usepackage{amssymb}
\let\Bbbk\relax
\usepackage{newtxmath}
\usepackage{pifont} % for cross mark

\definecolor{darkergreen}{RGB}{0,100,0} % Darker green
\definecolor{darkerred}{RGB}{139,0,0} % Darker red
\definecolor{lightergreen}{RGB}{34,139,34} % Lighter green, still subdued
\definecolor{lighterred}{RGB}{205,92,92} % Lighter red, still subdued
% faThumbsOUp means the unfilled one
% faThumbsUp mean the solid one
\newcommand{\positive}{\textcolor{lightergreen}{\faThumbsOUp}} % Positive symbol
\newcommand{\negative}{\textcolor{lighterred}{\faThumbsODown}} % Negative symbol
\definecolor{mydarkgreen}{HTML}{2E5E17}

\SetCommentSty{textit}
\newcommand\mygreen[1]{\textcolor{mydarkgreen}{\textit{#1}}}
\SetCommentSty{mygreen}

% chapter and section ref
\usepackage{cleveref}
\crefformat{section}{\S#2#1#3} % see manual of cleveref, section 8.2.1
\crefformat{subsection}{\S#2#1#3}
\crefformat{subsubsection}{\S#2#1#3}

% Chart Makecell
\usepackage{makecell}

% Listing
\usepackage{xcolor} % Required for custom colors
\usepackage{listings} % Required for listings
% Define custom colors

% listing color and font
% \usepackage{wasysym}
\usepackage{pifont}% http://ctan.org/pkg/pifont
\definecolor{dmlgreen}    {RGB}{51,  160,  44}
\definecolor{dmlblue}     {RGB}{31,  120, 180}
\definecolor{dmlred}      {RGB}{202,   0,  32}
\definecolor{brown}       {RGB}{139,  69,  19}

\definecolor{deepblue}{rgb}{0,0,0.5}
\definecolor{deepred}{rgb}{0.6,0,0}
\definecolor{deepgreen}{rgb}{0,0.5,0}
\definecolor{mauve}{rgb}{0.58,0,0.82}
\definecolor{light-gray}{gray}{0.96}
\definecolor{aliceblue}{rgb}{0.94, 0.97, 1.0}
\definecolor{OliveGreen}{rgb}{0.33, 0.42, 0.18}

\DeclareMathOperator*{\argmax}{arg\,max}
\DeclareMathOperator*{\argmin}{arg\,min}

\usepackage[utf8]{inputenc}
\DeclareFixedFont{\ttb}{T1}{txtt}{bx}{n}{8} % for bold
\DeclareFixedFont{\ttm}{T1}{txtt}{m}{n}{8}  % for normal

\newcommand{\updated}[1]{\noindent{\color{teal} {#1}}}
\newcommand{\zhanda}[1]{\noindent{\color{orange} {\bf Zhanda: } {#1}}}

\usepackage{xargs}                      % Use more than one optional parameter in a new commands
\usepackage[colorinlistoftodos,prependcaption,textsize=small]{todonotes}
\newcommandx{\note}[2][1=]{\todo[linecolor=OliveGreen,backgroundcolor=OliveGreen!25,bordercolor=OliveGreen,#1]{#2}}
\newcommandx{\unsure}[2][1=]{\todo[linecolor=red,backgroundcolor=red!25,bordercolor=red,#1]{#2}}
\newcommandx{\improve}[2][1=]{\todo[linecolor=orange,backgroundcolor=orange!25,bordercolor=orange,#1]{#2}}
\newcommandx{\change}[2][1=]{\todo[linecolor=Plum,backgroundcolor=Plum!25,bordercolor=Plum,#1]{#2}}
\newcommandx{\fix}[2][1=]{\todo[linecolor=blue,backgroundcolor=blue!25,bordercolor=blue,#1]{#2}}

% Spacing between figure and its caption.
% \usepackage{titlesec}
\setlength{\abovecaptionskip}{5pt plus 2pt minus 2pt}
\setlength{\belowcaptionskip}{-6pt}
% Spacing of section titles. Format:
% {left (don't touch)}{above}{below}
% \titlespacing\section{0pt}{5pt plus 4pt minus 2pt}{5pt plus 2pt minus 2pt}
% \titlespacing\subsection{0pt}{5pt plus 4pt minus 2pt}{5pt plus 2pt minus 2pt}
% \titlespacing\subsection{0pt}{5pt plus 4pt minus 2pt}{5pt plus 2pt minus 2pt}
% \titlespacing\subsubsection{0pt}{5pt plus 4pt minus 2pt}{5pt plus 2pt minus 2pt}

% Python style for highlighting
\newcommand\pythonstyle{\lstset{
  language=Python,
  backgroundcolor=\color{white}, %%%%%%%
  basicstyle=\ttfamily\footnotesize,
  otherkeywords={self},
  keywordstyle=\footnotesize\color{deepblue},
  emph={pipeline,unroll,tile},          
  emphstyle=\footnotesize\color{deepred},    
  stringstyle=\color{deepgreen},
  commentstyle=\color{red},  %%%%%%%%
  frame=tb,                         
  showstringspaces=false,
}}
\lstnewenvironment{env_python}[2]
{
\lstset{
  language=Python,
  backgroundcolor=\color{light-gray},
  otherkeywords={self},        
  emph={to,tile,reuse_at,reorder,unroll,split,vectorize,quantize},          
  emphstyle=\footnotesize\color{dmlred},  
  belowcaptionskip=0.7\baselineskip,
  aboveskip=-2mm,
  belowskip=5mm,
  showstringspaces=false,
  columns=flexible,
  escapechar=@,
  basicstyle={\linespread{0.99}\fontencoding{T1}\footnotesize\fontfamily{lmtt}\fontseries{m}\selectfont},
  numbers={left},
  xleftmargin={2em},%
  framesep=5pt, % margin top and down
  framexleftmargin={1em},
  numberstyle=\tiny\color{gray},
  keywordstyle=\color{blue},
  commentstyle=\color{deepgreen},
  stringstyle=\color{mauve},
  breaklines=true,
  breakatwhitespace=true,
  tabsize=3,
  frame=tb, 
  caption= #1, 
  label= #2
}%lstset
  \vspace{\baselineskip}
}%lstnewenvironment
{}

\lstset{
  language=Python,
  backgroundcolor=\color{aliceblue},
  % otherkeywords={self},
  emph={},          
  emphstyle=\small\color{dmlred},  
  belowcaptionskip=0.7\baselineskip,
  aboveskip=0mm,
  belowskip=3mm,
  showstringspaces=false,
  columns=flexible,
  basicstyle={\linespread{1.1}\fontencoding{T1}\scriptsize\fontfamily{lmtt}\fontseries{m}\selectfont},
  numbers={left},
  xleftmargin={2em},%
  numberstyle=\tiny\color{gray},
  keywordstyle=\color{blue},
  commentstyle=\color{deepgreen},
  stringstyle=\color{mauve},
  breaklines=true,
  breakatwhitespace=true,
  tabsize=3,
  morecomment=[f][\lstbg{red!20}]-,
  morecomment=[f][\lstbg{green!20}]+,
  morecomment=[f][\textit]{@@},
 }

\newcommand{\lstbg}[3][0pt]{{\fboxsep#1\colorbox{#2}{\strut #3}}}
\lstdefinelanguage{diff}{
  backgroundcolor=\color{aliceblue},           
  emph={},          
  emphstyle=\small\color{dmlred},  
  belowcaptionskip=0.7\baselineskip,
  aboveskip=0mm,
  belowskip=3mm,
  showstringspaces=false,
  columns=flexible,
  basicstyle={\linespread{1.1}\fontencoding{T1}\scriptsize\fontfamily{lmtt}\fontseries{m}\selectfont},
  numbers={left},
  xleftmargin={2em},%
  breaklines=true,
  breakatwhitespace=true,
  tabsize=3,
  morecomment=[f][\lstbg{red!20}]-,
  morecomment=[f][\lstbg{green!20}]+,
  morecomment=[f][\textit]{@@},
}

\makeatletter
\let\old@lstKV@SwitchCases\lstKV@SwitchCases
\def\lstKV@SwitchCases#1#2#3{}
\makeatother
\usepackage{lstlinebgrd}
\makeatletter
\let\lstKV@SwitchCases\old@lstKV@SwitchCases

\lst@Key{numbers}{none}{%
    \def\lst@PlaceNumber{\lst@linebgrd}%
    \lstKV@SwitchCases{#1}%
    {none:\\%
     left:\def\lst@PlaceNumber{\llap{\normalfont
                \lst@numberstyle{\thelstnumber}\kern\lst@numbersep}\lst@linebgrd}\\%
     right:\def\lst@PlaceNumber{\rlap{\normalfont
                \kern\linewidth \kern\lst@numbersep
                \lst@numberstyle{\thelstnumber}}\lst@linebgrd}%
    }{\PackageError{Listings}{Numbers #1 unknown}\@ehc}}
\makeatother

% ========================================================
\newcommand{\myparagraph}[1]{\vspace{1pt}\noindent\textbf{{#1}.\;}}
\newcommand{\cmark}{\ding{51}} % pifont checkmark
\newcommand{\xmark}{\ding{55}} % pifont cross
\newcommand{\filledone}{\ding{182}\xspace}
\newcommand{\filledtwo}{\ding{183}\xspace}
\newcommand{\filledthree}{\ding{184}\xspace}
\newcommand{\filledfour}{\ding{185}\xspace}
\newcommand{\filledfive}{\ding{186}\xspace}
\newcommand{\mynewpage}{\newpage\null\newpage}

\definecolor{customcolor}{HTML}{f9ae78}
\definecolor{customcolor2}{HTML}{3d5c6f}

\newcommand{\update}[1]{\noindent{\color{cyan} {\bf } {#1}}}

\usepackage{xcolor}
\usepackage{listings}
\lstdefinestyle{mystyle}{
%     backgroundcolor=\color{cyan!20}, % Light blue background
    basicstyle=\ttfamily\small,      % Small font
%     keywordstyle=\color{blue},       % Blue for keywords like "cd"
    numbers=left,                    % Line numbers on the left
      numberstyle=\color{gray},   % Small gray numbers
%     frame=single,                    % Border around the listing
%     rulecolor=\color{gray!50},       % Gray border color
%     breaklines=true                   % Allow line breaks
    % xleftmargin=0pt,                 % No left margin
    % framerule=0pt,
    breakautoindent=false,
    breakindent=10pt,
    escapechar=|,
}
\lstset{style=mystyle}

\usepackage{tikz}
\usetikzlibrary{patterns}
\usepackage[precision=2, unit=mm]{lengthconvert}
\newcommand{\dummyfigure}[3]{%
    \begin{tikzpicture}
        \draw [
            thick, 
            pattern=north east lines, 
            pattern color=black!10, 
            inner sep=0,
        ] 
        (0,0) rectangle (#1 - \pgflinewidth, #2 - \pgflinewidth) 
        node[
            pos=0.5, 
            align=center, 
            font=\small, 
            text width=#1,
        ] {
            #3 \\[2mm] 
            \texttt{\tiny{Figure size: (\Convert{#1}, \Convert{#2})}}
        };
    \end{tikzpicture}%
}